\ifdefined\gkver\else\def\gkver{student}\fi
\documentclass[\gkver]{gaokaozhenti}
\begin{document}

\chapter{1985年全国}

\begin{enumerate}
\item
%% number: 1
%% paperName: 1985年全国高考第 1 题 3 分
%% typeId: 3
%% type: 填空题
%% chapter: 磁场
%% point: 安培力
%% method: 无
%% score: 3
%% degree: 650
%% duplicateId: 0
%% body:
在两条平行长直导线中通以相同方向的电流时，它们之间的作用为互相\tkanswer{吸引}；通以相反方向的电流时，则互相\tkanswer{排斥}。

%% answer: 吸引，排斥
\jdanswer{
吸引，排斥
}

%% memo:
\memoanswer{
\onepicture[w=4cm]{figs/01_an.png}

在两条平行长直导线中通以相同方向的电流时，如图所示，由右手螺旋定则判断，左边电流在导线右侧产生的磁场方向向里，由左手定则判断，右边电流的受力方向向左；同样的方法可以判断左边电流的受力方向向右。所以它们之间的作用为互相吸引；同理可判断在两条平行长直导线中通以相反方向的电流时，则互相排斥。

评分标准：答对一个给 1 分，答对两个给 3 分。
}

\item
%% number: 2
%% paperName: 1985年全国高考第 2 题 3 分
%% typeId: 3
%% type: 填空题
%% chapter: 原子和原子核
%% point: 玻尔模型
%% method: 无
%% score: 3
%% degree: 750
%% duplicateId: 0
%% body:
有一群处在量子数 $n=3$ 的激发态中的氢原子，在它们的发光过程中发出的光谱线共有\tkanswer{3}条。

%% answer: 3
\jdanswer{
3
}

%% memo:
\memoanswer{
跃迁的可能情况如下：从 $n=3$ 到 $n=2$，从 $n=3$ 到 $n=1$，从 $n=2$ 到 $n=1$，共 3 条。
}

\item
%% number: 3
%% paperName: 1985年全国高考第 3 题 3 分
%% typeId: 3
%% type: 填空题
%% chapter: 机械振动
%% point: 等效单摆
%% method: 无
%% score: 3
%% degree: 750
%% duplicateId: 0
%% body:
图中所示为一双线摆，它是在一水平天花板上用两根等长细绳悬挂一小球而构成的，绳的质量可以忽略，设图中的 $l$ 和 $\alpha$ 为已知量。当小球垂直于纸面作简谐振动时，周期为\tkanswer{$T=2\pi\sqrt{\frac{l\sin\alpha}{g}}$}。
\onepicture[align=right, w=4cm]{\includesvg{figs/03}}

%% answer: 2\pi\sqrt{\frac{l\sin\alpha}{g}}
\jdanswer{
$T=2\pi\sqrt{\frac{l\sin\alpha}{g}}$
}

%% memo:
\memoanswer{
双线摆相当于单摆，摆长为 $l\sin\alpha$，代入单摆的周期公式 $T=2\pi\sqrt{\frac{l}{g}}$，得 $T=2\pi\sqrt{\frac{l\sin\alpha}{g}}$。
}

\item
%% number: 4
%% paperName: 1985年全国高考第 4 题 3 分
%% typeId: 3
%% type: 填空题
%% chapter: 原子和原子核
%% point: 天然放射性
%% method: 无
%% score: 3
%% degree: 850
%% duplicateId: 0
%% body:
$\ce{^{238}_{92}U}$ 衰变成 $\ce{^{234}_{90}Th}$ 时释放出\tkanswer{$\alpha$}粒子；$\ce{^{234}_{90}Th}$ 衰变成 $\ce{^{234}_{91}Pa}$ 时释放出\tkanswer{$\beta$}粒子；如果衰变时产生的新核处于激发态，将会辐射出\tkanswer{$\gamma$ 光子（或 $\gamma$ 射线）}。

%% answer: \alpha，\beta，\gamma 光子（或 \gamma 射线）
\jdanswer{
$\alpha$，$\beta$，$\gamma$ 光子（或 $\gamma$ 射线）
}

%% memo:
\memoanswer{
原子核的衰变方程为 $\ce{^{238}_{92}U -> ^{234}_{90}Th + ^{4}_{2}He}$，$\ce{^{234}_{90}Th -> ^{234}_{91}Pa + ^{0}_{-1}e}$，由衰变方程可知，释放的粒子分别为 $\alpha$ 粒子和 $\beta$ 粒子；若新核处于激发态，还会向外辐射 $\gamma$ 光子。
}

\item
%% number: 5
%% paperName: 1985年全国高考第 5 题 3 分
%% typeId: 3
%% type: 填空题
%% chapter: 光学
%% point: 光的折射
%% method: 无
%% score: 3
%% degree: 750
%% duplicateId: 0
%% body:
附图为一摄谱仪的示意图。来自光源的光经过狭缝 S 和透镜 L 后，成为平行光射在三棱镜 P 上。为了在一照相底片 MN 上拍摄下清晰的光谱，在 P 与 MN 之间须放置一个\tkanswer{凸透镜}。黄、红、绿三种颜色光谱线在照相底片上从 M 端到 N 端的次序为\tkanswer{红黄绿}。
\onepicture[align=right, w=8cm]{figs/05.png}

%% answer: 凸透镜，红黄绿
\jdanswer{
凸透镜，红黄绿
}

%% memo:
\memoanswer{
为了使光线更集中，放凸透镜。又因为光线的频率越大，折射率越大，根据光谱线的频率从小到大的顺序为红黄绿，故从 M 端到 N 端的次序为红黄绿。
}

\item
%% number: 6
%% paperName: 1985年全国高考第 6 题 3 分
%% typeId: 3
%% type: 填空题
%% chapter: 电路
%% point: 电流强度
%% method: 无
%% score: 3
%% degree: 550
%% duplicateId: 0
%% body:
设 a 和 b 为长度相同的两段均匀纯铜丝，它们的截面积之比为 $S_{a}:S_{b}=1:2$。在每一铜丝两端都加以相同电压 $U$，这时两铜丝中自由电子的定向运动速度之比为 $v_{a}:v_{b}=$\tkanswer{$1:1$}。

%% answer: 1:1
\jdanswer{
$1:1$
}

%% memo:
\memoanswer{
由电阻定律 $R=\rho\frac{l}{S}$ 知 $\frac{R_{a}}{R_{b}}=\frac{S_{b}}{S_{a}}=\frac{2}{1}$，由欧姆定律得 $I=\frac{U}{R}$，则 $\frac{I_{a}}{I_{b}}=\frac{R_{b}}{R_{a}}=\frac{1}{2}$，又电流的微观表达式为 $I=neSv$，得 $v=\frac{I}{neS}$，故 $\frac{v_{a}}{v_{b}}=\frac{I_{a}}{I_{b}}\cdot\frac{S_{b}}{S_{a}}=\frac{1}{2}\times\frac{2}{1}=\frac{1}{1}$。
}

\item
%% number: 7
%% paperName: 1985年全国高考第 7 题 3 分
%% typeId: 3
%% type: 填空题
%% chapter: 功和能
%% point: 动能定理
%% method: 建立模型
%% score: 3
%% degree: 350
%% duplicateId: 0
%% body:
某地强风的风速约为 $v=20\Ums$，设空气密度为 $\rho=1.3\Ukgmc$。如果把通过横截面积为 $S=20\Umq$ 的风的动能全部转化为电能，则利用上述已知量计算电功率的公式应为 $P=$\tkanswer{$\frac{1}{2}\rho Sv^{3}$}，大小约为\tkanswer{$10^{5}$}$\UW$（取一位有效数字）。

%% answer: \frac{1}{2}\rho Sv^{3}，10^{5}
\jdanswer{
$\frac{1}{2}\rho Sv^{3}$，$10^{5}$
}

%% memo:
\memoanswer{
单位时间内风的动能转化为电能为 $P=\frac{1}{2}mv^{2}$，又 $m=\rho Sv$，解得 $P=\frac{1}{2}\rho Sv^{3}$，代入数据得 $P=1\times10^{5}\UW$。
}

\item
%% number: 8
%% paperName: 1985年全国高考第 8 题 3 分
%% typeId: 1
%% type: 单选题
%% chapter: 曲线运动
%% point: 圆周运动的应用
%% method: 无
%% score: 3
%% degree: 650
%% duplicateId: 0
%% body:
一圆盘可绕一通过圆盘中心 O 且垂直于盘面的竖直轴转动，在圆盘上放置一木块，当圆盘匀角速转动时，木块随圆盘一起运动。那么\xzanswer{B}
\onepicture[w=3.2cm]{figs/08.png}
\fivechoices[answer=B]
{木块受到圆盘对它的摩擦力，方向背离圆盘中心}
{木块受到圆盘对它的摩擦力，方向指向圆盘中心}
{因为木块随圆盘一起运动，所以木块受到圆盘对它的摩擦力，方向与木块的运动方向相同}
{因为摩擦力总是阻碍物体运动，所以木块所受圆盘对它的摩擦力的方向与木块的运动方向相反}
{因为二者是相对静止的，圆盘与木块之间无摩擦力}

%% answer: B
%% memo:
\memoanswer{
木块随圆盘一起做匀速圆周运动的向心力由摩擦力提供，所以木块受到圆盘对它的摩擦力，方向指向圆盘中心，故 B 正确。
}

\item
%% number: 9
%% paperName: 1985年全国高考第 9 题 3 分
%% typeId: 1
%% type: 单选题
%% chapter: 气体性质
%% point: 物体的内能
%% method: 无
%% score: 3
%% degree: 750
%% duplicateId: 0
%% body:
有一定质量的气体，其温度由 $T_{1}$ 升高到 $T_{2}$，在这过程中，\xzanswer{D}
\fourchoices[answer=D]
{如果气体体积膨胀并因而对外界作功，则分子的平均平动能可能会减少}
{如果气体体积保持不变，则分子的平均平动能可能不变}
{只有当气体体积被压缩并因而外界对气体作功时，分子的平均平动能才会增加}
{不管气体的体积如何变化，分子的平均平动能总是增加的}

%% answer: D
%% memo:
\memoanswer{
温度是平均动能的标志，温度升高，气体分子的平均动能增大，故 D 正确。
}

\item
%% number: 10
%% paperName: 1985年全国高考第 10 题 3 分
%% typeId: 1
%% type: 单选题
%% chapter: 磁场
%% point: 带电粒子在磁场中的运动
%% method: 无
%% score: 3
%% degree: 550
%% duplicateId: 0
%% body:
ab 是一弯管，其中心线是半径为 $R$ 的一段圆弧。将它置于一给定的匀强磁场中，磁场方向垂直于圆弧所在平面（即纸面）并且指向纸外：有一束粒子对准 a 端射入弯管，粒子有不同的质量、不同的速度，但都是一价正离子。\xzanswer{C}
\onepicture[w=6.5cm]{figs/10.png}
\fourchoices[answer=C]
{只有速度大小一定的粒子可以沿中心线通过弯管}
{只有质量大小一定的粒子可以沿中心线通过弯管}
{只有动量大小一定的粒子可以沿中心线通过弯管}
{只有能量大小一定的粒子可以沿中心线通过弯管}

%% answer: C
%% memo:
\memoanswer{
带电粒子在磁场中做圆周运动的向心力由洛伦兹力提供，有 $qvB=m\frac{v^{2}}{R}$，得 $R=\frac{mv}{qB}$，因 $B$、$q$ 一定，要求 $R$ 一定，必须 $mv$ 一定，故 C 正确。
}

\item
%% number: 11
%% paperName: 1985年全国高考第 11 题 3 分
%% typeId: 2
%% type: 多选题
%% chapter: 光学
%% point: 光电效应
%% method: 无
%% score: 3
%% degree: 550
%% duplicateId: 0
%% body:
已知一光电管的阴极的极限频率为 $\nu_{0}$，现将频率 $\nu$ 大于 $\nu_{0}$ 的光照射在阴极上。\xzanswer{ACD}
\onepicture[w=3.5cm]{figs/11.png}
\fourchoices[answer=ACD]
{照射在阴极上的光的强度愈大，单位时间内产生的光电子数目也愈多}
{加在 AK 间的正向电压愈大，通过光电管的光电流饱和值也愈大}
{为了阻止光电子到达 A，必须在 AK 间加一足够高的反向电压}
{阴极材料的逸出功等于 $h\nu_{0}$}

%% answer: ACD
%% memo:
\memoanswer{
照射在阴极上的光的强度愈大，则单位时间内光子数目愈多，所以单位时间内产生的光电子数目也愈多，故 A 正确；通过光电管的光电流饱和值由照射在阴极上的光的强度决定，加在 AK 间的正向电压不能改变，故 B 错误；在 AK 间加一足够高的反向电压，对光电子减速，可以使光电子不能到达 A，故 C 正确；阴极材料的逸出功，由定义知 $W_{0}=h\nu_{0}$，故 D 正确。
}

\item
%% number: 12
%% paperName: 1985年全国高考第 12 题 3 分
%% typeId: 2
%% type: 多选题
%% chapter: 电场
%% point: 电场强度
%% method: 无
%% score: 3
%% degree: 550
%% duplicateId: 0
%% body:
电场强度 $E$ 的定义式为 $E=\frac{F}{q}$。\xzanswer{BD}
\fourchoices[answer=BD]
{这定义式只适用于点电荷产生的电场}
{上式中，$F$ 是放入电场中的电荷所受的力，$q$ 是放入电场中的电荷的电量}
{上式中，$F$ 是放入电场中的电荷所受的力，$q$ 是产生电场的电荷的电量}
{在库仑定律的表达式 $F=k\frac{q_{1}q_{2}}{r^{2}}$ 中，$k\frac{q_{2}}{r^{2}}$ 是点电荷 $q_{2}$ 产生的电场在点电荷 $q_{1}$ 处的场强大小；而 $k\frac{q_{1}}{r^{2}}$ 是点电荷 $q_{1}$ 产生的电场在点电荷 $q_{2}$ 处的场强大小}

%% answer: BD
%% memo:
\memoanswer{
由电场强度的定义可以判断 BD 正确。
}

\item
%% number: 13
%% paperName: 1985年全国高考第 13 题 3 分
%% typeId: 1
%% type: 单选题
%% chapter: 电磁感应
%% point: 楞次定律
%% method: 无
%% score: 3
%% degree: 650
%% duplicateId: 0
%% body:
一均匀的扁平条形磁铁与一圆形线圈同在一平面内，磁铁中央与圆心 O 重合。为了在磁铁开始运动时在线圈中得到一方向如图所示的感生电流 $i$，磁铁的运动方式应为\xzanswer{A}
\onepicture[w=3.6cm]{figs/13.png}
\sixchoices[answer=A]
{N 极向纸内，S 极向纸外，使磁铁绕 O 点转动}
{N 极向纸外，S 极向纸内，使磁铁绕 O 点转动}
{使磁铁沿垂直于线圈平面的方向向纸外作平动}
{使磁铁沿垂直于线圈平面的方向向纸内作平动}
{使磁铁在线圈平面内绕 O 点沿顺时针方向转动}
{使磁铁在线圈平面内绕 O 点沿逆时针方向转动}

%% answer: A
%% memo:
\memoanswer{
N 极向纸内相当于圆形线圈的上段向纸外切割磁感线运动，用右手定则判断感应电流方向如题图所示，故 A 正确 B 错误；磁铁垂直纸面向外或向内运动，圆形线圈的上段与下段产生的感应电流方向相反，互相抵消，故 CD 错误；磁铁在纸面内绕 O 点转动，圆形线圈内磁通量没有变化，不产生感应电流，故 EF 错误。
}

\item
%% number: 14
%% paperName: 1985年全国高考第 14 题 3 分
%% typeId: 1
%% type: 单选题
%% chapter: 曲线运动
%% point: 竖直平面圆周运动
%% method: 无
%% score: 3
%% degree: 350
%% duplicateId: 0
%% body:
如图，一细绳的上端固定在天花板上靠近墙壁的 O 点，下端拴一小球。L 点是小球下垂时的平衡位置。Q 点代表一固定在墙上的细长钉子，位于 OL 直线上。N 点在 Q 点正上方，且 QN $=$ QL。M 点与 Q 点等高。现将小球从竖直位置（保持绳绷直）拉开到与 N 等高的 P 点，释放后任其向 L 摆动。运动过程中空气阻力可忽略不计。小球到达 L 后，因细绳被长钉挡住，将开始沿以 Q 为中心的圆弧继续运动。在这以后，\xzanswer{E}
\onepicture[w=4.2cm]{figs/14.png}
\fivechoices[answer=E]
{小球向右摆到 M 点，然后就摆回来}
{小球向右摆到 M 和 N 之间圆弧上某点处，然后竖直下落}
{小球沿圆弧摆到 N 点，然后竖直下落}
{小球将绕 Q 点旋转，直到细绳完全缠绕在钉上为止}
{关于小球的运动情况，以上说法都不正确}

%% answer: E
%% memo:
\memoanswer{
由机械能守恒知，小球若沿圆弧摆到 N 点，则速度为 0，由圆周运动的规律知，小球若沿圆弧摆到 N 点，则在 N 点的速度至少为 $\sqrt{gR}$（其中 $R$ 表示以 Q 为中心的圆弧运动的半径），二者互相矛盾，所以小球沿圆弧不能摆到 N 点，故 C 错误；小球向右摆到 M 点时，速度不为 0，所以不能摆回来，故 A 错误；由内轨脱离模型知，小球向右摆到 M 和 N 之间圆弧上某点处时，速度不为 0，不能竖直下落，且脱离轨道后，小球做斜抛运动，故 B、D 错误。综上可知 E 正确。
}

\item
%% number: 15
%% paperName: 1985年全国高考第 15 题 2 分
%% typeId: 4
%% type: 实验题
%% chapter: 直线运动
%% point: 游标卡尺和螺旋测微仪
%% method: 无
%% score: 2
%% degree: 850
%% duplicateId: 0
%% body:
下图中表示用一零点准确的游标卡尺（主尺上每一小格等于 0.1 厘米）测量一圆柱体的直径。从放大了的插图中读出的测量结果为直径 $D=$\tkanswer{$2.23$}厘米。
\onepicture[align=right, w=6.5cm]{figs/15.png}

%% answer: 2.23
\jdanswer{
2.23
}

%% memo:
\memoanswer{
主尺读 2.2 厘米，游标尺的第 3 个格与主尺的某格对齐，应为 $0.1\times3=0.3\Umm$，故实际读数为 $2.23\Ucm$。

评分标准：本小题 2 分。答对的，给 2 分；答作 2.220 的同样给分。
}

\item
%% number: 16
%% paperName: 1985年全国高考第 16 题 5 分
%% typeId: 4
%% type: 实验题
%% chapter: 电路
%% point: 多用表的使用
%% method: 无
%% score: 5
%% degree: 650
%% duplicateId: 0
%% body:
下图为一可供使用的万用表。S 为选择开关，Q 为欧姆挡调零旋钮。现在要用它检验两个电阻的阻值（图中未画出电阻），已知阻值分别为 $R_{1}=60\UO$ 和 $R_{2}=470\UkO$。下面提出了在测量过程中一系列可能的操作。请你选出能尽可能准确地测定各阻值和符合于万用表安全使用规则的各项操作，并且将它们按合理顺序填写在后面的横线上空白处。

（A）旋动 S 使其尖端对准欧姆挡 $\times1$ k（即 OHMS$\times1$ K）

（B）旋动 S 使其尖端对准欧姆挡 $\times100$

（C）旋动 S 使其尖端对准欧姆挡 $\times10$

（D）旋动 S 使其尖端对准欧姆挡 $\times1$

（E）旋动 S 使其尖端对准 V1000

（F）将两表笔分别接到 $R_{1}$ 的两端，读出 $R_{1}$ 的阻值。随后即断开

（G）将两表笔分别接到 $R_{2}$ 的两端，读出 $R_{2}$ 的阻值。随后即断开

（H）两表笔短接，调节 Q 使表针对准欧姆挡刻度盘上的 0，随后即断开

所选操作及其顺序为（用字母代号填写）：\tkanswer[6cm]{DHFAHGE}。

（操作总数应视实际需要而定）
\onepicture[align=right, w=7cm]{figs/16.png}

%% answer: DHFAHGE
\jdanswer{
DHFAHGE
}

%% memo:
\memoanswer{
本题解题关键：
\begin{steps}
	\item
	每测一次进行一次欧姆挡调零；
	\item
	测量时要使表针转到一半以上，即正确选择挡位，如测量低电阻用欧姆挡 $\times1$，测量高电阻用欧姆挡 $\times1$k；
	\item
	测量完成后，要旋动 S 使其尖端对准 V1000；要旋动 S 使其尖端对准 V1000 的作用是保护万用电表。
\end{steps}

评分标准：本小题 5 分，DHF 和 AHG 各占 2 分，末尾的 E 占 1 分。答案为 AHGDHFE 的，同样给分。将 DHF 答为 CHF 的，扣 1 分；将 AHG 答为 BHG 的，扣 1 分。答案中缺一个 H（欧姆挡调 0）的，扣 1 分；两个 H 都缺的，扣 3 分。
}

\item
%% number: 17
%% paperName: 1985年全国高考第 17 题 7 分
%% typeId: 4
%% type: 实验题
%% chapter: 电路
%% point: 电路实验
%% method: 无
%% score: 7
%% degree: 450
%% duplicateId: 0
%% body:
有一改装的安培表 A$_{1}$ 需要与一标准安培表 A$_{2}$ 进行核对，采用如图 \ref{1985全国17a} 所示的电路，其中 $E$ 为电源，$R_{0}$ 为一限流电阻，$R$ 为一可变电阻，K 为电键。限流电阻能够限制住电路中的最大电流，使之不超出安培表的量程过多，从而对安培表起保护作用。实验中已有的器件及其规格如下：

蓄电池 $E$（电动势 $6\UV$，内阻约为 $0.3\UO$），

改装的安培表 A$_{1}$（量程 $0\sim0.6\UA$，内阻约为 $0.1\UO$），标准安培表 A$_{2}$（量程 $0\sim0.6\sim3\UA$，内阻不超过 $0.04\UO$）。

实验中备用的电阻器及其规格如下：

（A）固定电阻（阻值 $8\UO$，额定电流 $2\UA$）

（B）固定电阻（阻值 $15\UO$，额定电流 $2\UA$）

（C）滑动变阻器（阻值范围 $0\sim20\UO$，额定电流 $2\UA$）

（D）滑动变阻器（阻值范围 $0\sim200\UO$，额定电流 $2\UA$）
\begin{enumerate}
	\item
	已知两个表的刻度盘上都将量程均分为 6 大格，要求从 $0.1\UA$ 起对每条刻线一一进行核对。为此，从备用的电阻器中，$R_{0}$ 应选用\tkanswer{A}，$R$ 应选用\tkanswer{D}。（用字母代号填写）
	\item
	图 \ref{1985全国17b} 所示为电路中所需器件，在图上画出连线，将所示器件按电路图连接成实验电路。
\end{enumerate}
\twopicture[numA=1, numB=2, labelA=1985全国17a, labelB=1985全国17b, labelprefix={图}, wA=6cm, wB=7cm, align=right]
{figs/17a.png}
{figs/17b.png}

%% answer: （1）A，D；（2）电路图连接如图所示
\jdanswer{
\begin{enumerate}
	\item
	A，D

	\item
	电路图连接如图所示
\end{enumerate}
}

%% memo:
\memoanswer{
电路中的最大电流 $I\approx\frac{6}{8}\UA\approx0.7\UA$，用阻值 $15\UO$ 的固定电阻，电路中的最大电流 $I\approx\frac{6}{15}\UA\approx0.4\UA$，后者不能满足核对 $0.6\UA$ 电流的要求，应选用阻值 $8\UO$ 的固定电阻，故应选择 A。

选择滑动变阻器：如用 C，最小电流为 $I\approx\frac{6}{28}\UA\approx0.2\UA$，不能满足对 $0.1\UA$ 电流的校准，故应选择 D。

电路图连接如图所示。

\onepicture[w=8cm]{figs/17_an.png}

评分标准：本小题共 7 分。元件选择占 3 分，只答对一个的，给 1 分。电路连线占 4 分，连法只要正确，虽与参考答案不同，不影响得分。将两个安培表并联的，给 0 分；将一个安培表并联在电源两端的，给 0 分；将电源短路或未连入电路中的，给 0 分；未将变阻器连入电路中或连成不能调的，给 0 分；安培表正负端连接错误的，扣 2 分；标准安培表量程选 $0\sim3\UA$ 的，扣 2 分。
}

\item
%% number: 18
%% paperName: 1985年全国高考第 18 题 6 分
%% typeId: 6
%% type: 计算题
%% chapter: 功和能
%% point: 交通工具启动
%% method: 无
%% score: 6
%% degree: 550
%% duplicateId: 0
%% body:
在一水平长直轨道上，一动力车牵引一质量为 $m=5000\Ukg$ 的小车厢以 $v_{0}=36\Ukmh$ 的速度匀速行驶。这时动力车对该车厢的输出功率为 $P=15000\UW$。如果使车厢与动力车脱开，车厢将滑行多长的距离而停止？

%% answer: s=167 m
\jdanswer{
$s=167\Um$
}

%% memo:
\memoanswer{
因为车厢原来作匀速运动，所以车厢所受牵引力的大小等于阻力的大小。因此，车厢所受的阻力 $f$ 为
$f=\frac{P}{v_{0}}$（a）

脱开动力车后，车厢在阻力作用下滑行的距离若为 $s$，根据功能关系可知
$fs=\frac{1}{2}mv_{0}^{2}$（b）

由以上两式得
$s=\frac{mv_{0}^{3}}{2P}$（c）

代入数值得
$s=\frac{5000\times10^{3}}{2\times15000}\approx167\Um$。

评分标准：本题 6 分。列出（a）式的，给 2 分；列出（b）式的，给 2 分；得出（c）式的，再给 1 分；最后答数正确的，再给 1 分。
}

\item
%% number: 19
%% paperName: 1985年全国高考第 19 题 7 分
%% typeId: 6
%% type: 计算题
%% chapter: 气体性质
%% point: 气体连接体
%% method: 无
%% score: 7
%% degree: 450
%% duplicateId: 0
%% body:
内径均匀的 U 形管中装入水银。两管中水银面与管口的距离均为 $l=10.0\Ucm$。大气压强为 $p_{0}=75.8\UcmHg$ 时，将右侧管口密封，然后从左侧管口处将一活塞缓慢向下推入管中，直到左右两侧水银面高度差达 $h=6.0\Ucm$ 时为止。求活塞在管内移过的距离。
\onepicture[align=right, w=4.5cm]{figs/19.png}

%% answer: x=6.4 cm
\jdanswer{
$x=6.4\Ucm$
}

%% memo:
\memoanswer{
压强单位使用 cmHg，令 $p_{1}$、$p_{2}$ 分别代表管中左、右两侧气体的压强，$\rho$ 代表水银密度，则有
$p_{1}-p_{2}=\rho gh$（a）

令 $x$ 代表活塞移动的距离，由玻意耳定律可知

对左侧气体：$p_{0}l=p_{1}\left(l+\frac{h}{2}-x\right)$（b）

对右侧气体：$p_{0}l=p_{2}\left(l-\frac{h}{2}\right)$（c）

由以上三式可解出 $x$，得
$x=\frac{\rho gh\left(l^{2}-\frac{h^{2}}{4}\right)+p_{0}lh}{p_{0}l+\rho gh\left(l-\frac{h}{2}\right)}$（d）

将 $p_{0}=75.8\UcmHg$，$\rho gh=6.0\UcmHg$，$l=10.0\Ucm$ 和 $h=6.0\Ucm$ 代入得 $x=6.4\Ucm$。

评分标准：本题 7 分。列出（a）式的，给 1 分；列出（b）式的，给 2 分；列出（c）式的，给 2 分；得出（d）式的，再给 1 分；最后答数正确的，再给 1 分。
}

\item
%% number: 20
%% paperName: 1985年全国高考第 20 题 7 分
%% typeId: 6
%% type: 计算题
%% chapter: 机械波
%% point: 波的多解性
%% method: 无
%% score: 7
%% degree: 550
%% duplicateId: 0
%% body:
一列横波在 $x$ 轴线上传播着，在 $t_{1}=0$ 和 $t_{2}=0.005\Us$ 时的波形曲线如图所示。
\begin{enumerate}
	\item
	由图中读出波的振幅和波长；
	\item
	设周期大于（$t_{2}-t_{1}$），如果波向右传播，波速多大？如果波向左传播，波速又是多大？
	\item
	设周期小于（$t_{2}-t_{1}$）并且波速为 $6000\Ums$，求波的传播方向。
\end{enumerate}
\onepicture[align=right, w=11cm]{figs/20.png}

%% answer: （1）A=0.2 m，\lambda=8 m；（2）v_{右}=400 m/s，v_{左}=1200 m/s；（3）波向左传播
\jdanswer{
\begin{enumerate}
	\item
	振幅 $A=0.2\Um$，波长 $\lambda=8\Um$。

	\item
	$v_{\text{右}}=400\Ums$，$v_{\text{左}}=1200\Ums$。

	\item
	波向左传播。
\end{enumerate}
}

%% memo:
\memoanswer{
（1）由图知振幅 $A=0.2\Um$，波长 $\lambda=8\Um$。

（2）当 $(t_{2}-t_{1})$ 小于一个周期时，波的传播距离小于一个波长。如果波向右传播，则传播距离由图上可看出为 $2\Um$。由此得
$v_{\text{右}}=\frac{s_{\text{右}}}{t_{2}-t_{1}}=\frac{2}{0.005}\Ums=400\Ums$。

如果波向左传播，由图可知传播距离为 $6\Um$，由此得
$v_{\text{左}}=\frac{s_{\text{左}}}{t_{2}-t_{1}}=\frac{6}{0.005}\Ums=1200\Ums$。

（3）解法一：波的传播距离为 $s=v(t_{2}-t_{1})=6000\times0.005\Um=30\Um$，比三个波长多出 $6\Um$。由图可知波是向左传播的。

解法二：如果波向右传播，传播的最短距离为 $2\Um=\frac{1}{4}\lambda$，由波的周期性得 $2\Um=n\lambda+\frac{1}{4}\lambda$，则波速为
$v=\frac{n\lambda+\frac{1}{4}\lambda}{\Delta t}=\frac{8n+2}{0.005}\Ums=1600n+400\Ums$，$n$ 取自然数。

当 $n=3$ 时，$v=5200\Ums$；

当 $n=4$ 时，$v=6800\Ums$；

波速无法取 $6000\Ums$，故波不可能向右传播。

如果波向左传播，传播的最短距离为 $6\Um=\frac{3}{4}\lambda$，由波的周期性得 $6\Um=n\lambda+\frac{3}{4}\lambda$，则波速为
$v=\frac{n\lambda+\frac{3}{4}\lambda}{\Delta t}=\frac{8n+6}{0.005}\Ums=1600n+1200\Ums$，$n$ 取自然数。

当 $n=3$ 时，$v=6000\Ums$，波速可取得 $6000\Ums$，故波向左传播。

评分标准：全题 7 分。（1）2 分，（2）4 分，（3）1 分。（1）中，振幅和波长各占 1 分。（2）中，两个传播速度各占 2 分。（3）中，只有最后答案而无算式或解释说明的，不给分。
}

\item
%% number: 21
%% paperName: 1985年全国高考第 21 题 12 分
%% typeId: 6
%% type: 计算题
%% chapter: 电场
%% point: 带电粒子的圆周运动
%% method: 无
%% score: 12
%% degree: 450
%% duplicateId: 0
%% body:
如图所示，一条长为 $l$ 的细线，上端固定，下端拴一质量为 $m$ 的带电小球。将它置于一匀强电场中，电场强度大小为 $E$，方向是水平的。已知当细线离开竖直位置的偏角为 $\alpha$ 时，小球处于平衡。
\begin{enumerate}
	\item
	小球带何种电荷？求出小球所带电量；
	\item
	如果使细线的偏角由 $\alpha$ 增大到 $\varphi$，然后将小球由静止开始释放，则 $\varphi$ 应为多大，才能使在细线到达竖直位置时小球的速度刚好为零？
\end{enumerate}
\onepicture[align=right, w=4.5cm]{\includesvg{figs/21}}

%% answer: （1）正电，q=\frac{mg\tan\alpha}{E}；（2）\varphi=2\alpha
\jdanswer{
\begin{enumerate}
	\item
	小球带正电，$q=\frac{mg\tan\alpha}{E}$。

	\item
	$\varphi=2\alpha$。
\end{enumerate}
}

%% memo:
\memoanswer{
（1）由小球所受电场力的方向与场强方向相同，可知小球带正电。

小球受三个力作用：重力 $mg$，线拉力 $T$，电场力 $qE$，平衡时
$T\cos\alpha=mg$（a）

$T\sin\alpha=qE$（b）

由以上两式可得电量为
$q=\frac{mg\tan\alpha}{E}$（c）

\onepicture[w=3.5cm]{\includesvg{figs/21_an1}}

（2）当小球由初始位置 1 运动到竖直位置 2 时，线对它的拉力不做功，且已知原来的动能和后来的动能都等于 0，根据能量守恒可知，重力势能的减少量应等于电势能的增加量，因为重力势能的减少量为
$\Delta E_{\mathrm{Gp}}=mgl(1-\cos\varphi)$（d）

电势能的增加量 $\Delta E_{p}=qEl\sin\varphi$（e）

所以得
$mgl(1-\cos\varphi)=qEl\sin\varphi$（f）

化简得
$\frac{1-\cos\varphi}{\sin\varphi}=\frac{qE}{mg}$，则 $\tan\frac{\varphi}{2}=\frac{qE}{mg}=\tan\alpha$，即 $\varphi=2\alpha$（g）

故细线与竖直方向的夹角为 $2\alpha$ 时，满足要求。

\onepicture[w=3.5cm]{\includesvg{figs/21_an2}}

评分标准：全题 12 分。（1）4 分，（2）8 分。（1）中，正确判断出电荷正负的，给 1 分；（a）、（b）两式列对一个的，给 1 分，两个都列对的，给 2 分；得出（c）式的，再给 1 分；没列（a）、（b）式直接写出（c）式的，也给 3 分。（2）中，正确列出重力势能的变化（或重力做功）即（d）式的，给 2 分；正确列出电势能的变化（或电场力做功）即（e）式的，给 2 分；列出（f）式的，给 3 分；得出最后结果（g）式的，再给 1 分。
}

\item
%% number: 22
%% paperName: 1985年全国高考第 22 题 12 分
%% typeId: 6
%% type: 计算题
%% chapter: 电场
%% point: 带电粒子在交变电场中的运动
%% method: 无
%% score: 12
%% degree: 350
%% duplicateId: 0
%% body:
图 \ref{1985全国22a} 中 A 和 B 表示在真空中相距为 $d$ 的两平行金属板。加上电压后，它们之间的电场可视为匀强电场。图 \ref{1985全国22b} 表示一周期性的交变电压的波形，横坐标代表时间 $t$，纵坐标代表电压 $U$。从 $t=0$ 开始，电压为一给定值 $U_{0}$，经过半个周期，突然变为 $-U_{0}$；再过半个周期，又突然变为 $U_{0}$；$\ldots\ldots$如此周期性地交替变化。

在 $t=0$ 时，将上述交变电压 $U$ 加在 A、B 两板上，使开始时 A 板电势比 B 板高，这时在紧靠 B 板处有一初速为零的电子（质量为 $m$，电量为 $e$）在电场作用下开始运动。要想使这电子到达 A 板时具有最大的动能，则所加交变电压的频率最大不能超过多少？
\twopicture[numA=1, numB=2, labelA=1985全国22a, labelB=1985全国22b, labelprefix={图}, wA=2.4cm, wB=8cm, align=right]
{figs/22a.png}
{figs/22b.png}

%% answer: \nu=\sqrt{\frac{eU_{0}}{8md^{2}}}
\jdanswer{
频率不能超过 $\nu=\sqrt{\frac{eU_{0}}{8md^{2}}}$。
}

%% memo:
\memoanswer{
开始 $t=0$ 时，因 A 板电势比 B 板高，而电子又紧靠 B 板处，所以电子将在电场力作用下向 A 板运动。在交变电压的头半个周期内，电压不变，电子做匀加速直线运动，其动能不断增大。如果频率很高，即周期很短，在电子尚未到达 A 板之前交变电压已过了半个周期开始反向，则电子将沿原方向开始做匀减速直线运动。再过半个周期后，其动能减小到零。接着又变为匀加速运动，半个周期后，又做匀减速运动，……最后到达 A 板。

在匀减速运动过程中，电子动能要减少。因此，要想电子到达 A 板时具有最大的动能，在电压的大小给定了的条件下，必须使电子从 B 到 A 的过程中始终做加速运动。这就是说，要使交变电压的半周期不小于电子从 B 板处一直加速运动到 A 板处所需的时间，即频率不能大于某一值。

在电场力的作用下，电子的加速度 $a$ 为
$a=\frac{eU_{0}}{md}$（a）

其中 $e$ 和 $m$ 分别为电子的电量大小和质量，令 $t$ 表示电子从 B 一直加速运动到 A 所需的时间，则
$d=\frac{1}{2}at^{2}$ 或 $t=\sqrt{\frac{2d}{a}}$（b）

令 $T$ 表示交变电压的周期，$\nu$ 表示频率，根据以上的分析，它们应满足以下的要求：
$t\leq\frac{T}{2}$（c）

即 $\nu\leq\frac{1}{2t}$（d）

由（a）、（b）、（d）三式可解得
$\nu=\sqrt{\frac{eU_{0}}{8md^{2}}}$（e）

即频率不能超过 $\sqrt{\frac{eU_{0}}{8md^{2}}}$。

评分标准：本题 12 分。列出（a）式的，给 2 分；列出（b）式的，给 2 分；这一部分共占 4 分。能按题意得出条件（c）式或（d）式的，给 5 分；得出最后结果（e）式的，再给 3 分。这一部分共占 8 分。
}

\item
%% number: 23
%% paperName: 1985年全国高考第 23 题 10 分
%% typeId: 6
%% type: 计算题
%% chapter: 运动和力
%% point: 动量和能量
%% method: 无
%% score: 10
%% degree: 350
%% duplicateId: 0
%% body:
一质量为 $m_{1}$ 的入射粒子与一质量为 $m_{2}$ 的静止粒子发生正碰，已知机械能在碰撞过程中有损失，实验中测出了碰撞后第二个粒子的速度为 $v_{2}$，求第一个粒子原来速度 $v_{0}$ 的值的可能范围。

%% answer: \frac{m_{1}+m_{2}}{2m_{1}}v_{2}<v_{0}<\frac{m_{1}+m_{2}}{m_{1}}v_{2}
\jdanswer{
$\frac{m_{1}+m_{2}}{2m_{1}}v_{2}<v_{0}<\frac{m_{1}+m_{2}}{m_{1}}v_{2}$
}

%% memo:
\memoanswer{
粒子碰撞前后，由动量守恒得
$m_{1}v_{0}=m_{1}v_{1}+m_{2}v_{2}$①

其中 $v_{1}$ 是碰撞后第一个粒子的速度，由于机械能有损失，则
$\frac{1}{2}m_{1}v_{0}^{2}>\frac{1}{2}m_{1}v_{1}^{2}+\frac{1}{2}m_{2}v_{2}^{2}$②

为了消去 $v_{1}$，可将以上两式改写为
$m_{1}v_{0}^{2}-m_{1}v_{1}^{2}>m_{2}v_{2}^{2}$，$m_{1}v_{0}-m_{1}v_{1}=m_{2}v_{2}$

将上两式左右两侧分别取商，得 $v_{0}+v_{1}>v_{2}$，或写作 $v_{1}>v_{2}-v_{0}$③

将③式代入①式，得
$v_{0}>\frac{m_{1}+m_{2}}{2m_{1}}v_{2}$④

上式给出了 $v_{0}$ 的下限，即 $\frac{m_{1}+m_{2}}{2m_{1}}v_{2}$，此值对应于机械能守恒的情形。

碰撞后，第二个粒子必定是沿第一个粒子碰撞前的速度方向运动，所以 $v_{0}\leq v_{2}$⑤，代入①式得
$v_{0}\leq\frac{m_{1}+m_{2}}{m_{1}}v_{2}$⑥

上式给出了 $v_{0}$ 的上限，即 $\frac{m_{1}+m_{2}}{m_{1}}v_{2}$，此值对应于碰撞后两个粒子合并为一而运动的情形。

综上，速度范围为 $\frac{m_{1}+m_{2}}{2m_{1}}v_{2}<v_{0}<\frac{m_{1}+m_{2}}{m_{1}}v_{2}$。
}

\end{enumerate}

\end{document}
